\addcontentsline{toc}{chapter}{Abstrak}
\chapter*{Abstrak}

\section*{Abstrak}

Dimensi metrik dan dimensi metrik sisi pada sebuah graf merupakan invarian fundamental dalam teori graf. Secara sederhana, dimensi metrik menyatakan banyaknya simpul paling sedikit yang diperlukan untuk membedakan seluruh simpul lain dalam graf, sedangkan dimensi metrik sisi menyatakan hal serupa tetapi untuk membedakan seluruh sisi. Invarian ini memiliki berbagai aplikasi, seperti navigasi jaringan, pelacakan sumber, robotika, dan sistem biologis.

Tugas akhir ini mengeksplorasi penggunaan \foreign{Satisfiability Modulo Theories} (SMT) sebagai pendekatan untuk menghitung dimensi metrik dan dimensi metrik sisi. Permasalahan diformulasikan secara deklaratif dalam SMT dan diselesaikan menggunakan \foreign{solver}. Representasi graf kanonik digunakan untuk menghapus redundansi akibat automorfisma, dan eksperimen dilakukan pada seluruh graf terhubung unik dengan 2 hingga 9 simpul, serta beberapa graf berukuran lebih besar.

\section*{Kata Kunci}

Dimensi Metrik, Dimensi Metrik Sisi, \foreign{Satisfiability Modulo Theories}, SMT \foreign{Solver}, \foreign{NP-hard}.

\cleardoublepage

\addcontentsline{toc}{chapter}{Abstract}
\chapter*{Abstract}

\section*{Abstract}

The metric dimension and the edge metric dimension of a graph are fundamental invariants in graph theory. Roughly speaking, the metric dimension is the minimum number of vertices needed to distinguish all other vertices in the graph, while the edge metric dimension does the same for the edges. These invariants have a range of applications, such as in network navigation, source tracking, robotics, and biological systems.

This thesis explores the use of Satisfiability Modulo Theories (SMT) as an approach to compute the metric dimension and the edge metric dimension. The problems are formulated declaratively in SMT and solved using a solver. Canonical graph representations are employed to remove redundancy caused by automorphisms, and experiments are conducted on all unique connected graphs with 2 to 9 vertices, as well as on some larger graphs.

\section*{Key words}

Metric Dimension, Edge Metric Dimension, Satisfiability Modulo Theories, SMT Solver, NP-hard.
